
% \yz{Considering we haven't used RLHF yet (?), it's probably better to call this section finetuning or instruction tuning.}
% \yc{I'd say the term `alignment' is okay to use here as well, as alignment covers both instruction learning and feedback/preference learning in the current context. }
% \alex{I vote for finetuning. Alignment is a bigger topic than SFT. Using it will set the wrong expectation which makes the audience to expect RLHF / safety and all other stuffs, which we do not have yet}
% In most cases, directly using pre-trained large language models does not align with\qj{are not in accord with} human values. The use of alignment techniques, such as supervised fine-tuning (SFT) and reinforcement learning from human feedback (RLHF), can significantly improve the model's ability to serve as a helpful and harmless AI assistant. \ww{Maybe "Instruction Tuning" is a better subsitution for "supervised fine-tuning" in here ? SFT is a term of a big scope, "Instruction Tuning" is a more specific technique and more suitable to be listed in parallel with RLHF.} In this section, we will delve into the details of how Yi models have been trained\qj{aligned} using SFT. \footnote{Due to time constraints, we have not included the RLHF step.}

Our finetuning method significantly emphasizes data quality over quantity.
Our approach does \textit{not} follow existing data-intensive approaches like FLAN~\citep{chung2022scaling} and UltraChat~\citep{ding2023enhancing}, which scales the SFT data to millions of entries but each of the entries may not been examined carefully because the scale is too large.
In contrast, our method aligns with the LIMA~\citep{zhou2023lima} and DEITA~\citep{liu2023makes} approach, which focus on data selection rather than scaling. 
With the scale being less than 10K, we are able to examine and optimize \textit{every single data point}. 
Below we discuss our data construction and training details. 


\subsection{Data Preprocessing}

\paragraph{Quality is All You Need}
Our finetuning dataset consists of less than 10K multi-turn instruction-response dialog pairs, with each and every one of the entry constructed and polished over multiple iterations and from user feedback.
We take this approach because  in our preliminary experiments, we observe that compared to the open-source data of several hundred thousand entries, the results from a smaller, manually annotated dataset are superior.
% These findings are in line with the spirit \jc{key points} of those reported by \ty{Rephrase:
These observations align with those reported in~\citet{llama2, zhou2023lima, team2023gemini}. 
% Based on these insights, we implemented a lightweight instruction fine-tuning strategy for our alignment stage, using fewer than 10,000 
% data entries in total. To ensure quality, we have conducted strict manual rewriting for both the prompts and responses
% \ww{Rephrase: Is "response" a better word than "answer" ?} 
% for all finetuning data.

We use the following techniques to improve prompt distribution selection, response formatting, and chain-of-thought formatting:
(1). for prompt distribution selection,
drawing inspiration from WizardLM\citep{xu2023wizardlm}, we develope compound instructions and progressively evolved them to increase their complexity. 
This approach has significantly reduced the size of SFT data 
in our experiments;
(2). for response formatting, we generally use a default style extended from LIMA\cite{zhou2023lima}. Overall, the responses are structured in an introduction-body-conclusion format where the body is usually a list of bullet point;
(3). for CoT data formatting, we have use a ``Step-Back'' pattern, inspired by \citet{zheng2023step}, by performing abstraction to formulate higher-level solutions before delving into reasoning about the original, more concrete questions.

We spend extra efforts on reducing hallucination and repetition:
% Beyond the basic quality requirements, we have also paid extra attention to other quality requirements of the conversational assistant regarding the data. 
(1). to reduce hallucinations, 
we examine and ensure that the knowledge in the responses is not contained within the model, and eliminate responses that might lead to memorization;
(2). to reduce repetition, we rewrite the repetitive turns of the responses that usually exist but may be overlooked in the finetuning data. 

% For another instance,
% multi-turn data introduced from earlier models exhibit repetition between turns and we have rewritten these repetitive turns to aid the model in breaking free from these repetitive patterns.

\paragraph{Diversity and Mixture}
To ensure the coverage of different capabilities, we have included a wide spectrum of open-source prompt,
% \ty{Rephrase:a wide range of open-source prompts}, 
encompassing areas such as question answering, creative writing, dialogue, reasoning, mathematics, coding, safety, bilingual capabilities, and others. 

To obtain a fine-grained control of different directions of capabilities, inspired by InsTag\citep{lu2023instag}, we develop a instruction tagging system. By designing a diversity-focused sampling algorithm, we carefully balanced the distribution of instructions across various tags. This approach ensures a diverse finetuning dataset, aiming to achieve enhanced cross-task robustness.

To achieve the optimal data ratio for balancing different directions of the capability, we use an approximate grid search to determine our data mixture. Motivated by \citet{dong2023abilities}, this process involved experimenting with \{1, 1/2, 1/4, 1/8, 1/16, 1/32, 1/64\} proportions for each ability. The search process was guided by validation results 
% \ww{validation results? ``benchmark results'' can mislead to the impression of "validate on test set".} 
and our in-house human evaluation sets.

\paragraph{ChatML Format} Beyond the focus on data quality and diversity, our observations revealed that the format of the data substantially influences the model's ultimate performance. To this end, we implemented the ChatML-style format~\citep{chatml}. This structured approach empowers the model to differentiate among various information types, such as system configurations, user inputs, and assistant responses.

\subsection{Training Method}
% \yc{maybe title as `Model Training' or `Our Training Method'. }

% Similar to our pretraining methodology, next-token prediction remains the principal training task for each Supervised Fine-Tuning (SFT) experiment. 
We use next-word prediction loss for finetuning, and only compute loss on the responses, but not system and user instructions. 
% We strategically apply loss masks to both the system instructions and user inputs. 
We use AdamW optimizer
% Optimization is carried out using AdamW, with following hyperparameters: 
with $\beta_1$ set to 0.9, $\beta_2$ set to 0.999, and $\epsilon$ set to $10^{-8}$. 
% We maintain a sequence length of 4096, 
We use a sequence length of 4096, alongside a batch size of 64.
% The training 300 steps, all executed under a constant learning rate of $1\times 10^{-5}$. 
We set training step to 300 with a constant $1\times 10^{-5}$ learning rate, a weight decay of 0.1, gradient clipping with a maximum threshold of 1.0, and NEFTune~\cite{jain2023neftune} with a noise scale of 45 for Yi-34B-Chat and 5 for Yi-6B-Chat.
% To avoid overfitting, we implement a weight decay of 0.1, enforcing gradient clipping with a maximum threshold of 1.0 and NEFTune~\cite{jain2023neftune} with a noise scale of 45 for Yi-34B-Chat and 5 for Yi-6B-Chat.